\chapter{Diseases of the Reproductive Organs}
\label{ch:reproductive}
The reproductive organs include the male reproductive system (testes, epididymis, vas deferens, seminal vesicles, prostate, and penis) and the female reproductive system (ovaries, fallopian tubes, uterus, cervix, vagina, and vulva). Diseases range from benign conditions to malignant neoplasms, hormonal disorders, infections, and structural abnormalities.

%-------------------------------------------------------------
\section{Female Reproductive Diseases}
%-------------------------------------------------------------

%-------------------------------------------------------------

%-------------------------------------------------------------

%-------------------------------------------------------------

%-------------------------------------------------------------

%-------------------------------------------------------------

\label{sec:Female Reproductive Diseases}
%-------------------------------------------------------------%-------------------------------------------------------------

\subsection{Atrophic Vaginitis (Genitourinary Syndrome of Menopause)}
\label{sec:Atrophic Vaginitis}
Atrophic vaginitis, now termed genitourinary syndrome of menopause (GSM), is a chronic progressive condition resulting from estrogen deficiency after menopause, affecting up to 50\% of postmenopausal women. The condition results from estrogen withdrawal leading to thinning of the vaginal epithelium, decreased vaginal secretions, loss of rugae, shortening and narrowing of the vaginal canal, and changes in the vulvar and lower urinary tract tissues. Clinical features include vaginal dryness, burning, itching, dyspareunia, and urinary symptoms (frequency, urgency, recurrent UTIs). Diagnosis is clinical: pale, thin, smooth vaginal mucosa with loss of rugae, elevated vaginal pH ($>$5.0), and maturation index showing predominance of parabasal cells on cytology. Management includes vaginal moisturizers and lubricants for mild symptoms, and low-dose vaginal estrogen therapy for moderate to severe symptoms. Prognosis is excellent with treatment, significantly improving quality of life and sexual function.

\subsection{Cervical Cancer}
\label{sec:Cervical Cancer}
Cervical cancer is the fourth most common cancer in women worldwide, with approximately 600,000 new cases and 340,000 deaths annually, disproportionately affecting low- and middle-income countries. Nearly all cases (99\%) are caused by persistent infection with high-risk human papillomavirus (HPV) types 16 and 18 (responsible for 70\% of cases). The condition results from progression from HPV infection to invasive cancer over 10--20 years. Clinical features include abnormal vaginal bleeding, postcoital bleeding, and pelvic pain. Cervical screening involves Pap smear (cytology) and/or HPV DNA testing, with abnormal results evaluated by colposcopy and biopsy. Pre-invasive lesions (cervical intraepithelial neoplasia, CIN) are treated by excisional procedures (LEEP, cold knife conization). Management of early-stage disease is radical hysterectomy with pelvic lymph node dissection or chemoradiation; advanced disease is treated with concurrent chemoradiation and adjuvant chemotherapy. Prognosis is excellent with early detection, and prevention through HPV vaccination and screening has dramatically reduced incidence.

\subsection{Endometriosis}
\label{sec:Endometriosis}
Endometriosis is defined by the presence of endometrial-like tissue outside the uterus, affecting approximately 10\% of reproductive-age women, with the most common sites being the ovaries (endometriomas---``chocolate cysts''), uterosacral ligaments, pelvic peritoneum, and rectovaginal septum. The pathogenesis is debated, with Sampson's theory of retrograde menstruation being the most widely accepted. Clinical features include chronic pelvic pain (dysmenorrhea, dyspareunia, chronic pelvic pain), infertility (present in 30--50\% of women with endometriosis), and cyclic symptoms. Diagnosis is typically made by laparoscopy with visualization of endometriotic implants and histological confirmation, with MRI useful for deep infiltrating endometriosis. Management includes NSAIDs for pain, combined oral contraceptives, progestins, GnRH agonists/antagonists with add-back therapy, and surgical excision/ablation of endometriotic lesions. Prognosis is variable, with recurrence possible after treatment.

\subsection{Ovarian Cancer}
\label{sec:Ovarian Cancer}
Ovarian cancer is the fifth leading cause of cancer-related death in women, with approximately 313,000 new cases and 207,000 deaths annually worldwide, often diagnosed at an advanced stage due to vague symptoms. Risk factors include BRCA1/2 mutations (lifetime risk 40--60\% and 10--30\%, respectively), family history, nulliparity, early menarche, and late menopause, with protective factors including oral contraceptive use and tubal ligation. The most common histological type is high-grade serous carcinoma (70\%), arising from the fallopian tube epithelium. The condition results from malignant transformation of ovarian or fallopian tube epithelium. Clinical features include abdominal bloating, pelvic/abdominal pain, urinary urgency, and early satiety. CA-125 is a serum tumor marker used for monitoring treatment response. Diagnosis is by imaging, surgical exploration, and histopathology. Management is optimal cytoreductive surgery followed by platinum-based chemotherapy (carboplatin + paclitaxel), with maintenance therapy using PARP inhibitors for BRCA-mutated tumors. Prognosis is generally poor due to late diagnosis.

\subsection{Pelvic Inflammatory Disease}
\label{sec:Pelvic Inflammatory Disease}
Pelvic inflammatory disease is an infection of the female upper genital tract, including the endometrium (endometritis), fallopian tubes (salpingitis), ovaries (oophoritis), and pelvic peritoneum. It is the most common and most serious complication of sexually transmitted infections, with \textit{Chlamydia trachomatis} and \textit{Neisseria gonorrhoeae} being the most common causative organisms. Risk factors include multiple sexual partners, prior PID, younger age, and lack of barrier contraception. The condition results from ascending infection from the lower genital tract. Clinical features include lower abdominal pain (often bilateral), cervical motion tenderness (``chandelier sign''), uterine tenderness, adnexal tenderness, and abnormal cervical/vaginal discharge. Diagnosis is clinical (CDC criteria). Management includes empiric antibiotics covering \textit{N.~gonorrhoeae} and \textit{C.~trachomatis}: ceftriaxone plus doxycycline plus metronidazole. Prognosis is good with prompt treatment, though complications include tubo-ovarian abscess, infertility, and ectopic pregnancy.

\subsection{Polycystic Ovary Syndrome}
\label{sec:Polycystic Ovary Syndrome}
Polycystic ovary syndrome is the most common endocrine disorder in women of reproductive age, affecting 6--12\%. The condition results from insulin resistance (present in 50--70\%), hyperinsulinemia, and relative hyperandrogenism. Clinical features include menstrual irregularity, infertility, hirsutism, acne, obesity (60--80\%), metabolic syndrome, type 2 diabetes mellitus, and increased cardiovascular risk. Diagnosis requires two of three Rotterdam criteria: oligo-ovulation or anovulation, clinical or biochemical hyperandrogenism, and polycystic ovarian morphology on ultrasound, with exclusion of other causes. Management is tailored to goals: combined oral contraceptives for menstrual regulation and hyperandrogenism, metformin or inositol for insulin resistance, letrozole for ovulation induction in infertility, and lifestyle modification. Prognosis is good with appropriate management.

\subsection{Uterine Fibroids}
\label{sec:Uterine Fibroids}
Uterine fibroids (leiomyomas) are the most common benign tumors of the female reproductive tract, affecting up to 80\% of women by age 50, with higher prevalence and symptomatic disease in African American women. The condition results from monoclonal smooth muscle tumors of the myometrium, driven by estrogen and progesterone. Clinical features depend on location (submucosal, intramural, subserosal, pedunculated) and may include abnormal uterine bleeding (menorrhagia leading to iron deficiency anemia), pelvic pain and pressure, urinary frequency, constipation, infertility, and pregnancy complications. Diagnosis is by pelvic ultrasound (usually sufficient) or MRI. Management depends on symptoms and reproductive goals: medical management (combined oral contraceptives, progestins, GnRH agonists, tranexamic acid), uterine artery embolization, myomectomy, and hysterectomy. Prognosis is excellent, with fibroids being benign.

\subsection{Vaginal Cancer}
\label{sec:Vaginal Cancer}
Vaginal cancer is the rarest gynecologic malignancy, with squamous cell carcinoma (70--80\%) being the most common type, often related to HPV infection and preceded by vaginal intraepithelial neoplasia (VAIN). Risk factors include prior cervical cancer/VIN, HPV infection, smoking, and immunosuppression. The condition results from malignant transformation of vaginal epithelium. Clinical features include vaginal bleeding, discharge, dyspareunia, and pelvic pain. Diagnosis is by pelvic examination with biopsy. Management depends on stage and location: radiation therapy (most common treatment) or local excision for small, superficial lesions. Prognosis depends on stage at diagnosis.

\subsection{Vulvar Cancer}
\label{sec:Vulvar Cancer}
Vulvar cancer is a rare malignancy accounting for approximately 5\% of gynecologic cancers, typically affecting postmenopausal women. Squamous cell carcinoma (80--90\%) is the most common histological type, with risk factors including HPV infection (particularly in younger women), lichen sclerosus, vulvar intraepithelial neoplasia (VIN), smoking, and immunosuppression. The condition results from malignant transformation of vulvar epithelium. Clinical features include a vulvar lump, itching, pain, bleeding, and ulceration. Diagnosis requires biopsy of any suspicious lesion. Management is wide local excision with adequate margins, with inguinal lymph node dissection for invasive tumors $>$1 mm. Prognosis is related to stage and lymph node status.

%-------------------------------------------------------------
\section{Male Reproductive Diseases}
%-------------------------------------------------------------

%-------------------------------------------------------------

%-------------------------------------------------------------

%-------------------------------------------------------------

%-------------------------------------------------------------

%-------------------------------------------------------------

\label{sec:Male Reproductive Diseases}
%-------------------------------------------------------------%-------------------------------------------------------------

\subsection{Balanitis and Balanoposthitis}
\label{sec:Balanitis and Balanoposthitis}
Balanitis is inflammation of the glans penis; balanoposthitis additionally involves the foreskin, and is common in uncircumcised men. The condition results from infectious (Candida albicans---most common, bacterial, STIs), irritant (soap, detergent), or dermatologic (lichen planus, psoriasis, contact dermatitis) factors, with diabetes mellitus predisposing to candida balanitis. Clinical features include erythema, edema, pruritus, discharge, dysuria, and tender inguinal lymphadenopathy. Diagnosis is clinical; swab for KOH preparation (candida), bacterial culture, and STI screening. Management includes improved hygiene, topical antifungal (clotrimazole, miconazole) for candida, topical antibiotic/steroid combinations, and addressing underlying causes. Prognosis is excellent with treatment.

\subsection{Benign Prostatic Hyperplasia}
\label{sec:Benign Prostatic Hyperplasia}
Benign prostatic hyperplasia is the most common benign neoplasm in men, characterized by proliferation of stromal and epithelial cells in the transition zone of the prostate, causing urethral obstruction, affecting over 50\% of men by age 60 and up to 90\% by age 85. The condition results from age-related hormonal changes driving prostatic growth. Clinical features include lower urinary tract symptoms (LUTS): urinary frequency, urgency, nocturia, weak stream, hesitancy, straining, incomplete emptying, and dribbling, with acute urinary retention possibly occurring. Diagnosis involves digital rectal examination (DRE) showing smooth, symmetrically enlarged prostate, American Urological Association Symptom Score (AUASS), PSA, urinalysis, post-void residual volume, and uroflowmetry. Management depends on severity: watchful waiting for mild symptoms, alpha-blockers (tamsulosin, alfuzosin) for smooth muscle relaxation, 5-alpha-reductase inhibitors (finasteride, dutasteride) for gland shrinkage in larger prostates, and surgical intervention (TURP, holmium laser enucleation) for refractory cases. Prognosis is excellent with appropriate management.

\subsection{Erectile Dysfunction}
\label{sec:Erectile Dysfunction}
Erectile dysfunction is the persistent inability to achieve or maintain an erection sufficient for satisfactory sexual performance, affecting approximately 50\% of men between ages 40 and 70 to varying degrees. Risk factors include age, cardiovascular disease, diabetes mellitus, hypertension, dyslipidemia, smoking, obesity, depression, and medications, with the most common cause being vasculogenic (arterial insufficiency or veno-occlusive dysfunction), often reflecting subclinical atherosclerosis. The condition results from impaired vascular, neurological, or endocrine function. Clinical features include difficulty achieving or maintaining erections. Evaluation includes sexual history, hormonal testing (testosterone, prolactin, TSH), and cardiovascular risk assessment. First-line management is PDE5 inhibitors (sildenafil, tadalafil, vardenafil, avanafil), with second-line options including vacuum erection devices and intracavernosal injection therapy, and third-line being penile prosthesis implantation. Prognosis is good with appropriate management.

\subsection{Male Infertility}
\label{sec:Male Infertility}
Male infertility is responsible for approximately 40--50\% of couple infertility. Causes include varicocele (most common correctable cause, present in 40\% of infertile men), obstructive azoospermia (congenital absence of the vas deferens---often associated with CFTR mutations, epididymitis, vasectomy), non-obstructive azoospermia (Klinefelter syndrome, Y-chromosome microdeletions, cryptorchidism, prior chemotherapy/radiation), hypogonadotropic hypogonadism (Kallmann syndrome, pituitary tumors), and idiopathic oligozoospermia. The condition results from impaired spermatogenesis or sperm transport. Clinical features include inability to conceive after one year of unprotected intercourse. Evaluation includes semen analysis (WHO 2021 reference values), hormonal evaluation (FSH, LH, testosterone), and scrotal ultrasound. Management depends on the cause: varicocelectomy, assisted reproductive techniques (IUI, IVF with ICSI), surgical sperm retrieval, and treatment of reversible causes. Prognosis varies by etiology.

\subsection{Penile Cancer}
\label{sec:Penile Cancer}
Penile cancer is rare in developed countries (0.5--1 per 100,000) but accounts for up to 10\% of male malignancies in parts of Africa, South America, and Asia where circumcision is less common and HPV prevalence is higher, with squamous cell carcinoma representing $>$95\% of cases. Risk factors include lack of neonatal circumcision, phimosis, poor hygiene, HPV infection (types 16, 18, 45), smoking, and chronic balanoposthitis, with penile intraepithelial neoplasia (PeIN) being the precursor. The condition results from malignant transformation of squamous epithelium. Clinical features include a painless ulcer, mass, or wart-like lesion on the glans, prepuce, or shaft. Diagnosis is by biopsy, with staging using the TNM system. Management of the primary tumor includes topical therapy for PeIN, laser ablation or wide local excision for superficial tumors, and partial or total penectomy for invasive tumors. Prognosis is excellent for organ-confined disease; 5-year survival drops to 50\% with nodal involvement.

\subsection{Penile Fracture}
\label{sec:Penile Fracture}
Penile fracture is a traumatic rupture of the tunica albuginea of the corpus cavernosum, occurring during erectile activity when the penis is forcibly bent, most commonly during vigorous sexual intercourse or masturbation. Clinical features include a sudden snapping sound (audible crack), immediate detumescence, severe pain, and rapid development of swelling and ecchymosis (the ``eggplant deformity''), with urethral injury occurring in 10--20\% of cases. The condition results from traumatic rupture of the tunica albuginea. Diagnosis is clinical. Management is immediate surgical repair: degloving of the penis, evacuation of hematoma, identification and primary closure of the tunical defect. Prognosis is excellent with early surgical repair (erectile function preserved in $>$90\%), while conservative management leads to high rates of erectile dysfunction.

\subsection{Peyronie's Disease}
\label{sec:Peyronie's Disease}
Peyronie's disease is an acquired penile curvature caused by the formation of fibrous plaques in the tunica albuginea of the corpora cavernosa, affecting 3--9\% of men. The condition results from repetitive microtrauma during sexual activity in genetically susceptible individuals, leading to aberrant wound healing and excessive collagen deposition. Clinical features in the chronic phase (after 6--12 months) include stable plaque formation, penile curvature (dorsal, ventral, or lateral), palpable plaque, and often painful erections, with erectile dysfunction coexisting in 30--50\% of cases. Diagnosis is by history and physical examination; ultrasound or MRI can quantify plaque size and calcification. Management for the acute phase includes conservative management with observation, oral therapy, and intralesional collagenase, while chronic stable disease may require traction therapy, vacuum devices, intralesional verapamil or interferon, or surgical correction. Prognosis is variable.

\subsection{Phimosis and Paraphimosis}
\label{sec:Phimosis and Paraphimosis}
Phimosis is the inability to retract the foreskin over the glans penis, a normal finding in uncircumcised boys before age 5 (physiologic phimosis), while pathologic phimosis results from recurrent balanoposthitis, poor hygiene, or lichen sclerosus et atrophicus, causing scarring and a contracted, fibrous preputial ring. Paraphimosis is a urologic emergency in which the retracted foreskin cannot be returned to its normal position, constricting the glans and leading to vascular compromise, edema, and pain, typically occurring after catheterization, sexual activity, or inadequate reduction following examination. The condition results from preputial abnormalities. Clinical features of phimosis include difficulty with urination, recurrent infections, and painful erections; paraphimosis presents with a constricted glans and pain. Diagnosis is clinical. Management of paraphimosis involves gentle manual compression of the edematous glans followed by reduction; elective circumcision is the definitive treatment for pathologic phimosis. Prognosis is excellent with appropriate management.

\subsection{Priapism}
\label{sec:Priapism}
Priapism is a persistent, often painful, penile erection lasting longer than 4 hours unrelated to sexual stimulation. Ischemic (low-flow, veno-occlusive) priapism accounts for $>$95\% of cases and is a urologic emergency requiring prompt intervention to prevent corporal fibrosis and permanent erectile dysfunction. Causes include sickle cell disease (most common cause in children), intracavernosal injections, psychotropic medications, illicit drugs, anticoagulant therapy, and hematologic malignancies. Non-ischemic (high-flow, arterial) priapism results from unregulated arterial inflow, typically from perineal trauma causing arteriocavernosal fistula. The condition results from dysregulation of penile blood flow. Clinical features include a rigid, tender penis with soft glans in ischemic type. Diagnosis involves history, physical examination, and cavernous blood gas analysis. Management of ischemic priapism follows a stepwise approach: corporal aspiration with irrigation, intracavernosal phenylephrine, and surgical shunting. Prognosis depends on promptness of intervention.

\subsection{Prostate Cancer}
\label{sec:Prostate Cancer}
Prostate cancer is the second most commonly diagnosed cancer and the fifth leading cause of cancer-related death in men worldwide. Risk factors include age (rare under 50, risk increases sharply after 65), race (highest incidence in African Americans), family history, and genetic factors (BRCA2 mutations), with most prostate cancers being adenocarcinomas arising in the peripheral zone. The condition results from malignant transformation of prostatic epithelial cells. Clinical features are often asymptomatic in early stages, with advanced disease causing urinary symptoms, hematuria, and bone pain from metastases. Screening involves prostate-specific antigen (PSA) testing and digital rectal examination. Diagnosis is by transrectal or transperineal MRI-guided biopsy, with Gleason score and ISUP grade group guiding risk stratification. Management of low-risk disease may be active surveillance; intermediate-risk disease is treated with radical prostatectomy or radiation therapy; high-risk and metastatic disease require multimodal therapy including androgen deprivation therapy (ADT), radiation, chemotherapy (docetaxel), and novel hormonal agents (abiraterone, enzalutamide). Prognosis is excellent for localized disease and variable for advanced disease.

\subsection{Testicular Cancer}
\label{sec:Testicular Cancer}
Testicular cancer is the most common solid tumor in young men (peak age 15--40) and is highly curable. The vast majority (95\%) are germ cell tumors (GCTs): seminoma (40\%) and non-seminomatous germ cell tumors (NSGCT, 60\%)---embryonal carcinoma, yolk sac tumor, choriocarcinoma, teratoma, and mixed germ cell tumors. Risk factors include cryptorchidism, prior testicular cancer, Klinefelter syndrome, and infertility. The condition results from malignant transformation of germ cells. Clinical features typically include a painless testicular mass, though some present with scrotal pain, hydrocele, or symptoms of metastatic disease. Diagnosis involves scrotal ultrasound (hypoechoic intratesticular mass), serum tumor markers (AFP, $\beta$-hCG, LDH), and CT abdomen/pelvis/chest for staging. Management of stage I seminoma involves surveillance or adjuvant carboplatin; advanced disease is treated with cisplatin-based chemotherapy (BEP---bleomycin, etoposide, cisplatin) with surgery for residual masses. Prognosis is excellent, with cure rates exceeding 95\% for seminoma and 80--90\% for NSGCT.

